\scanpage{25}

Exemple d'un groupe lisse résoluble \inserted{connexe} $G$ sur $\mathbb F_p$, \inserted{de rang unipotent nul \deleted{\unsure{réductif|scindé} \illegible{1}} \unsure{linéaire|trivial},} dont l'algèbre
de Lie n'a pas d'élément régulier.

Soit $(\alpha_i)_{i\in I}$ une famille \inserted{finie} de formes \inserted{lin. $\ne0$} sur $\deleted{\unsure{\mathbb F_p^2|\mathbb F_p}}\inserted{\mathfrak t=\mathbb F_p^2}$ telles que
$\bigcup\operatorname{Ker}\alpha_i=\deleted{\unsure{\mathbb F_p|\mathbb F_p^2}(\mathbb F_p^2)}\mathfrak t\deleted{\mathbb F_p^2}$. (On peut s'en tirer
avec $p+1$ formes non nulles, et $\bigcap\operatorname{Ker}\alpha_i=0$). \deleted{Soit}
Les $\alpha_i$ sont les caractères infinitésimaux associés
à des \unsure{homs|hom|hom.}

\[
\chi_i:T=G_m^2\longrightarrow G_m
\]

qu'on utilise pour former le produit semi-direct

\[
G=T.U\deleted{G_a^I}\qquad U=G_a^I
\]

avec $T$ opérant sur $G_a^I$ par les $\chi_i$.
Alors $T$ opère sur $\mathfrak g/\mathfrak t$ par les caractères
$\chi_i$, ce qui montre que $T$ est un Cartan
i.e. le rang unipotent de $G$ est nul. Le
centre de $G$ est donc égal au centre réductif, donc
égal au noyau de $T\longrightarrow\operatorname{Gl}(\mathfrak g)$\,; qui est
nul, \deleted{\unsure{donc par suite}} Enfin, un \deleted{\unsure{par suite}} élément
\deleted{régulier} de $\mathfrak t$ \deleted{serait dans \unsure{le|la}} étant dans
$\bigcup\operatorname{Ker}\alpha_i$, ne peut être \uline{régulier}. Or \uline{$\unsure{\forall|\text{tt}|\text{tout}}$}
\uline{tore maximal} \uline{de $G$} est \uline{conjugué de} $T$
\deleted{\unsure{[comme|comme}} \inserted{(par un élément unique de $U(\mathbb F_p)$]} donc
\deleted{tout} $\mathfrak g$ n'a pas d'él\textsuperscript{mt} régulier [car
il serait conjugué d'un élément de $\mathfrak t$].

% slanted
\margin{\unsure{Ici|Donc}, $T$ est \deleted{\unsure{égal|éplt}} un sous-groupe de type $(C)$, i.e. $\mathfrak t$ est une sous-algèbre de Cartan}

% archive sheet number: 24
